\section{Agenten-Personas des Turing Generators}
\label{app:persona_definitions}
Die KI-gestützten Verarbeitungsschritte des Turing Generators verwenden
aufgabenspezifische Agentenrollen in Kombination mit unterschiedlichen
Personas. Während eine Rolle die auszuführende Aufgabe festlegt, definiert
die Persona die Perspektive, aus der diese Aufgabe bearbeitet wird.
Die nachfolgenden Definitionen dokumentieren die im finalen, evaluierten
Verarbeitungspfad eingesetzten Persona-Profile. Historisch vorhandene
Persona-Definitionen aus früheren Implementierungsständen werden nicht
aufgenommen. Die dargestellten Definitionsdateien werden unverändert aus dem
Implementierungsstand des Turing Generators wiedergegeben.

Die Persona-Ergebnisse stellen Verarbeitungsergebnisse beziehungsweise
Entscheidungsgrundlagen dar und besitzen keine eigenständige Engineering
Authority. Übereinstimmung zwischen mehreren Personas wird entsprechend
nicht mit fachlicher Richtigkeit gleichgesetzt. Die vollständigen
Definitionen werden im Folgenden nach ihrem Einsatz innerhalb des
Verarbeitungsprozesses gruppiert.

\subsection{Interpretation von Engineering-Information}
\label{app:personas_interpretation}

Für die Interpretation von Engineering-Information werden drei
unterschiedliche Perspektiven verwendet. Sie unterscheiden zwischen einer
möglichst quellennahen Interpretation, einer Systems-Engineering-Perspektive
und einer gezielt auf Mehrdeutigkeiten und Unsicherheiten ausgerichteten
Prüfperspektive.

\subsubsection*{Literal Interpreter}
\lstinputlisting[

basicstyle=\ttfamily\footnotesize,
breaklines=true,
columns=fullflexible
]{Anhang/PersonaDefinitions/literal_interpreter.md}

\subsubsection*{Systems Engineering Interpreter}
\lstinputlisting[

basicstyle=\ttfamily\footnotesize,
breaklines=true,
columns=fullflexible
]{Anhang/PersonaDefinitions/systems_engineering_interpreter.md}

\subsubsection*{Skeptical Ambiguity Interpreter}
\lstinputlisting[

basicstyle=\ttfamily\footnotesize,
breaklines=true,
columns=fullflexible
]{Anhang/PersonaDefinitions/skeptical_ambiguity_interpreter.md}

\subsection{Modellplatzierung}
\label{app:personas_model_placement}
Für die Zuordnung bereits autorisierter Engineering-Information zur
Zielmodellstruktur werden drei Modellierungsperspektiven verwendet.
Sie betrachten die zulässige Platzierung aus Sicht der Modellierungsregeln,
der Systemarchitektur sowie einer konservativen Gegenprüfung. Abweichungen
zwischen den Vorschlägen bleiben erhalten und können im anschließenden
Human Model Placement Review bewertet werden.
\subsubsection*{SysML Profile Modeler}
\lstinputlisting[
basicstyle=\ttfamily\footnotesize,
breaklines=true,
columns=fullflexible
]{Anhang/PersonaDefinitions/sysml_profile_modeler.md}
\subsubsection*{System Architecture Modeler}
\lstinputlisting[
basicstyle=\ttfamily\footnotesize,
breaklines=true,
columns=fullflexible
]{Anhang/PersonaDefinitions/system_architecture_modeler.md}
\subsubsection*{Conservative Modeling Reviewer}
\lstinputlisting[
basicstyle=\ttfamily\footnotesize,
breaklines=true,
columns=fullflexible
]{Anhang/PersonaDefinitions/conservative_modeling_reviewer.md}

\subsection{Model Quality Refinement}
\label{app:persona_model_quality}
Nach der fachlichen Autorisierung und Modellplatzierung kann ein
zusätzlicher KI-gestützter Refinement-Schritt zur sprachlichen und
strukturellen Verbesserung des Modellinhalts eingesetzt werden. Im Gegensatz
zu den vorhergehenden Multi-Persona-Gruppen verwendet dieser Schritt eine
einzelne konservativ ausgerichtete Editing-Persona. Die vorgeschlagenen
Änderungen besitzen ebenfalls keine eigenständige Engineering Authority.
\subsubsection*{Model Quality Refiner}
\lstinputlisting[
basicstyle=\ttfamily\footnotesize,
breaklines=true,
columns=fullflexible
]{Anhang/PersonaDefinitions/model_quality_refiner.md}